\section{Sharp volume bounds at fixed centre, radial width, and correlation}
\label{sec:volume-v2}
The determinant of the canonical correlation matrix controls the steering
ellipsoid volume: $\operatorname{Vol}(\mathcal E)=(4\pi/3)|\det S|$.
The preceding centre/radial correlation theorem admits a distinct sharp
volume refinement. In particular, the state maximizing correlation at a
fixed radial width generally does not maximize volume at that width.

\begin{theorem}[Exact signed-determinant interval]\label{thm:volume-radial-v2}
For canonical data $(c\nn,0,S)$ with $\norm\nn=1$, $0\leq c\leq1$ and
$t=\norm{S^{\mathsf T}\nn}\in[0,1-c]$, the exact signed-determinant range is
\begin{equation}\label{eq:signed-volume-v2}
 -\frac{t((1+t)^2-c^2)}4\ \leq\ \det S\ \leq\
 \frac{t((1-t)^2-c^2)}4.
\end{equation}
Every value in this interval is attained. Consequently the sharp absolute
volume bound is
\begin{equation}\label{eq:volume-radial-v2}
 \boxed{\quad |\det S|\leq \frac{t((1+t)^2-c^2)}4.\quad}
\end{equation}
For $t>0$ the absolute maximum has a unique local-unitary orbit. In radial
axes it is \eqref{eq:interval-family-v2} with the $00$--$11$ coherence
$\sqrt{ad}$ and the $01$--$10$ coherence zero. Its rank is three if
$t<1-c$, two if $0<c<1$ and $t=1-c$, and one at $(c,t)=(0,1)$.
At $t=0$, every feasible state has determinant zero.
\end{theorem}

\begin{proof}
Rotate into the radial axes used in
Theorem~\ref{thm:canonical-interval-v2}. The state then has the form
\begin{equation}\label{eq:aligned-volume-v2}
 \sigma=\begin{pmatrix}
 a&0&z&w\\0&b&v&-z\\
 \overline z&\overline v&e&0\\
 \overline w&-\overline z&0&d
 \end{pmatrix},
\end{equation}
where $a,b,e,d$ have the definitions in that theorem. Its actual correlation
matrix is
\[
 S=\begin{pmatrix}
 2\Re(w+v)&2\Im(v-w)&4\Re z\\
 -2\Im(w+v)&2\Re(v-w)&-4\Im z\\
 0&0&t
 \end{pmatrix}.
\]
Thus, putting $W=|w|^2$, $J=|v|^2$, $Z=|z|^2$,
\begin{equation}\label{eq:volume-coordinates-v2}
 \det S=4t(J-W),\qquad
 \tau=t^2+8(W+J+2Z).
\end{equation}
The two principal minors give $W\leq ad$ and $J\leq be$.
Since $ad-be=t/4\geq0$, these give precisely
\eqref{eq:signed-volume-v2} and \eqref{eq:volume-radial-v2}.
Take $z=0$ and independently choose
$0\leq W\leq ad$, $0\leq J\leq be$ with real nonnegative coherences.
The two $2\times2$ blocks are positive, so every signed determinant in the
claimed interval is attained.

If $t>0$, then $ad>be$. Absolute equality forces $W=ad$ and $J=0$.
The saturated $00,11$ principal minor and positivity force its two rows to
be proportional, which forces $z=0$. A phase rotation removes the phase
of $w$. The state is consequently the indicated unique orbit. Its block
ranks give the stated ranks. For $t=0$, the last row of $S$ vanishes and
the determinant is identically zero, so no uniqueness is asserted.
\end{proof}

At the extreme radial width $t=1-c$, \eqref{eq:volume-radial-v2} reduces to
$|\det S|\leq(1-c)^2$, the established fixed-centre maximum
of Milne et al.~\cite{milnevolume}. The refinement here fixes the radial
width as well, and then also fixes the correlation strength.
At smaller radial widths, the volume maximizer usually
has rank three, whereas the correlation maximizer has rank two. There is a
sharp quantitative tradeoff between these two objectives.

\begin{theorem}[Sharp correlation--volume tradeoff]\label{thm:triple-volume-v2}
Fix any feasible $c,t$ and any
$\tau\in[t^2,1-c^2+2t^2]$. Among canonical states with these three prescribed
parameters, the exact greatest absolute determinant is
\begin{equation}\label{eq:triple-volume-v2}
 \boxed{\quad
 V_{\max}(c,t,\tau)=\frac t2\min\bigl\{
 \tau-t^2,\ 1+2t+2t^2-c^2-\tau\bigr\}.
 \quad}
\end{equation}
In particular the largest volume occurs at
$\tau=t^2+[(1+t)^2-c^2]/2$, while at the largest correlation
$\tau=1-c^2+2t^2$ the determinant magnitude is at most $t^2$.
Every value on the upper tradeoff curve is attained by an explicit positive
$X$ state, including all rank-deficient and zero-width cases.
\end{theorem}

\begin{proof}
Use \eqref{eq:aligned-volume-v2} and set
\[
 U=ad=\frac{(1+t)^2-c^2}{16},\quad
 V=be=\frac{(1-t)^2-c^2}{16},\quad
 s=\frac{\tau-t^2}{8}=W+J+2Z.
\]
The allowed correlation interval gives $0\leq s\leq U+V$.
The block matrices $A=\operatorname{diag}(a,b)$ and
$D=\operatorname{diag}(e,d)$ satisfy
$A\preceq a\Id$ and $D\preceq d\Id$.
Block positivity gives, for arbitrary $\xi,\eta\in\C^2$,
\[
 |\xi^\dagger C\eta|^2\leq(\xi^\dagger A\xi)(\eta^\dagger D\eta)
                    \leq ad\norm\xi^2\norm\eta^2,
 \qquad C=\begin{pmatrix}z&w\\v&-z\end{pmatrix}.
\]
Here the scalar Cauchy--Schwarz inequality follows by restricting the
positive block matrix to the two vectors $(\xi,0)$ and $(0,\eta)$; it holds when
either restricted diagonal entry is zero. Consequently
$\norm C_{\mathrm{op}}^2\leq U$, and in particular
\[
 W+Z\leq U,\qquad J+Z\leq U.
\]
No inverse, strict positivity assumption, or limiting argument is needed.
The elementary inequalities
\[
 |W-J|\leq W+J\leq s,
 \qquad
 |W-J|=2\max\{W,J\}+2Z-s\leq2U-s
\]
therefore give $|\det S|\leq4t\min\{s,2U-s\}$.
Substituting the definitions yields \eqref{eq:triple-volume-v2}.

To attain the bound, set $z=0$ and take real nonnegative $w,v$. For
$0\leq s\leq U$, choose $W=s$, $J=0$. For $U\leq s\leq U+V$, choose
$W=U$, $J=s-U$. In each case $W\leq U$ and $J\leq V$, so both independent
$2\times2$ blocks of the state are positive. The correlation is exactly
$t^2+8s=\tau$. In the first branch the determinant magnitude is $4ts$;
in the second it is $4t(2U-s)$, where $2U-s\geq U-V=t/4\geq0$.
This proves sharpness for every feasible triple.
\end{proof}


\paragraph{Complete equality classification for the tradeoff.}
For $t>0$, keep $U=ad$, $V=be$, $s=(\tau-t^2)/8$, so that $U>V$.
If $0\leq s\leq U$, every maximizer, after removing coherence phases by
local $z$ rotations, is one of
\[
 (z,w,v)=(0,\sqrt s,0),\qquad
 (z,w,v)=(0,0,\sqrt s)\quad\text{when }s\leq V.
\]
The two representatives coincide at $s=0$. For $s>0$ they have opposite
signed determinant and hence belong to distinct local-unitary orbits.
If $U<s\leq U+V$, there is exactly one maximizing orbit, represented by
\[
 (z,w,v)=(0,\sqrt U,\sqrt{s-U}).
\]
To verify necessity in the first range, equality in
$|W-J|\leq W+J+2Z$ forces $Z=0$ and $\min\{W,J\}=0$.
In the second range, equality forces $\max\{W,J\}+Z=U$.
The second row actually satisfies $J+Z\leq bd<ad=U$, so the maximum must
be $W$. Applying positivity to the explicit vector
$(d,0,-\overline z,-\overline w)$ gives
\[
 0\leq d\,[ad-W-Z]-\frac t2Z.
\]
Since $t>0$, equality $W+Z=U$ forces $Z=0$ and then $W=U$.
The displayed phase removal uses two independent local phase rotations,
so it preserves all physical parameters and the signed determinant.
At $t=0$ every feasible state has zero volume and is a maximizer.

When $t<1-c$ and $t>0$, the first, negative-chirality branch has rank four
for $s<U$, rank three at $s=U$, rank three for $U<s<U+V$, and rank two
at $s=U+V$. The secondary positive-chirality branch has rank four for
$s<V$ and rank three at $s=V$. Zero diagonal entries at $t=1-c$ reduce
these ranks by the same block-counting rule used above.

The alignment of the two qubits used in these written proofs is realized by
local qubit unitaries. Signed-determinant invariance can be verified directly,
without assuming a correspondence with every real orthogonal matrix.
Using $\sigma_i\sigma_j=\delta_{ij}\Id+i\epsilon_{ijk}\sigma_k$, expansion
of a canonical density matrix gives
\[
 16\Tr(\sigma^3)=1+3c^2+3\tau-6\det S.
\]
Here $c^2=2\Tr(\sigma_A^2)-1$ and
$\tau=4\Tr(\sigma^2)-1-c^2$. A local unitary preserves each of these
three traces and keeps Bob maximally mixed. It therefore preserves
$\det S$ itself. This proves the asserted distinction between the two
signed orbits, and explains why the formulas describe chirality as well
as volume.

\paragraph{Formalization boundary.}
The finite-matrix determinant identity, both signed determinant bounds,
the rank-safe row quadratic-form tests, and the three-parameter upper
bound are checked in \path{MilneVolumeGeometry.lean}; the principal
statement is \path{alignedState_triple_volume_bound}.
The two independent positive coherence families and exact attained maximum
are checked in \path{MilneVolumeSharpness.lean}, principally by
\path{alignedTripleVolumeSpectrum_isGreatest} and
\path{volumeWitness_attains}.

The final results are stated for arbitrary density matrices in
\path{MilneVolumePhysical.lean}, principally by
\path{physicalTripleVolumeSpectrum_isGreatest} and
\path{physicalRadialVolumeSpectrum_isGreatest}. Their domain imposes the
actual Pauli conditions $\aa=c\mathbf e_z$, $\bb=0$, and
$\norm{S^{\mathsf T}\mathbf e_z}^2=t^2$, and imposes no form on the other
matrix entries. The axis $\mathbf e_z$ is a fixed marked axis, retained also
when $c=0$. The checked Bob-only normalization in
\path{MilneBobRadialNormalForm.lean} puts every such density matrix into
radial coordinates while preserving that marked Alice axis.
\path{MilneVolumeUnitary.lean} verifies the cubic-trace identity and the
signed-determinant invariance under actual local unitaries. Thus no
positivity surrogate, assumed row bound, assumed unitary invariance, or
rank assumption is supplied to the final physical theorems.
The additional orbit equality classification, complete signed-interval
attainability, and detailed rank counts in this section are written proofs;
they are not claimed as separately formalized Lean theorems.
